% =====================================================================
%  RELATÓRIO TÉCNICO — Projeto Prático: Sensor + Atuador com Arduino/ESP32
%  Introdução ao Hardware · Unoesc Chapecó · Prof. Jacson Luiz Matte
%
%  COMO USAR NO OVERLEAF
%  1. overleaf.com > New Project > Upload Project (envie este .tex) — ou
%     Blank Project e cole o conteúdo em main.tex.
%  2. Menu > Compiler: pdfLaTeX.
%  3. Coloque as imagens (diagrama, fotos) na pasta do projeto e ajuste os \includegraphics.
%  4. Substitua tudo que está entre << >>.
% =====================================================================
\documentclass[12pt,a4paper]{article}

\usepackage[utf8]{inputenc}
\usepackage[T1]{fontenc}
\usepackage[brazil]{babel}
\usepackage[margin=2.5cm]{geometry}
\usepackage{graphicx}
\usepackage{booktabs}
\usepackage{xcolor}
\usepackage{listings}
\usepackage{hyperref}
\usepackage{float}

% ---- identidade visual ----
\definecolor{unoescazul}{HTML}{1E4679}
\definecolor{unoescverde}{HTML}{46B652}
\definecolor{codefundo}{HTML}{F4F6F9}
\hypersetup{colorlinks=true, linkcolor=unoescazul, urlcolor=unoescazul, citecolor=unoescazul}

% ---- estilo de código Arduino ----
\lstdefinestyle{arduino}{
  language=C++,
  backgroundcolor=\color{codefundo},
  basicstyle=\ttfamily\small,
  keywordstyle=\color{unoescazul}\bfseries,
  commentstyle=\color{gray}\itshape,
  stringstyle=\color{unoescverde},
  numbers=left, numberstyle=\tiny\color{gray}, numbersep=8pt,
  frame=single, framerule=0pt, rulecolor=\color{codefundo},
  breaklines=true, showstringspaces=false, tabsize=2,
  morekeywords={pinMode,digitalWrite,digitalRead,analogRead,analogWrite,delay,millis,
                Serial,begin,print,println,available,read,HIGH,LOW,INPUT,OUTPUT,INPUT_PULLUP,
                setup,loop,byte,boolean,String}
}
\lstset{style=arduino}

\begin{document}

% =====================================================================
%  CAPA
% =====================================================================
\begin{titlepage}
  \centering
  {\large\bfseries\color{unoescazul} UNIVERSIDADE DO OESTE DE SANTA CATARINA\par}
  {\normalsize Campus Chapecó · Curso de Sistemas de Informação\par}
  {\normalsize 25830 · Introdução ao Hardware · 2026/2\par}
  \vspace{3cm}

  {\Large\bfseries\color{unoescazul} RELATÓRIO TÉCNICO\par}
  \vspace{0.6cm}
  {\LARGE\bfseries <<Nome do projeto>>\par}
  \vspace{0.4cm}
  {\large Projeto Prático — Sensor + Atuador com Arduino/ESP32\par}
  \vspace{0.3cm}
  {\normalsize Atividade Avaliativa III\par}
  \vspace{3cm}

  \begin{tabular}{l}
    <<Nome completo do integrante 1>> \\
    <<Nome completo do integrante 2>> \\
    <<Nome completo do integrante 3>> \\
    <<Nome completo do integrante 4>> \\
  \end{tabular}
  \vspace{2cm}

  {\normalsize Professor: Jacson Luiz Matte\par}
  \vfill
  {\normalsize Chapecó, <<mês>> de 2026\par}
\end{titlepage}

% =====================================================================
%  RESUMO + LINKS  (fica na primeira página do conteúdo)
% =====================================================================
\section*{Resumo}
<<Em até 8 linhas: qual problema o projeto resolve, quais sensores e atuadores foram usados,
qual placa (Arduino Uno ou ESP32) e qual foi o resultado. Escreva por último, quando tudo
estiver pronto.>>

\begin{center}
\fbox{\begin{minipage}{0.92\textwidth}
\textbf{Vídeo demonstrativo:} \url{https://<<link-do-video>>} \\[2pt]
\textbf{Código-fonte completo:} \url{https://<<link-do-repositorio-ou-tinkercad>>} \\[2pt]
\textbf{Circuito no Tinkercad (se houver):} \url{https://www.tinkercad.com/things/<<id>>}
\end{minipage}}
\end{center}

\tableofcontents
\newpage

% =====================================================================
\section{Introdução e problema}
<<Que situação do dia a dia o projeto ataca? Quem seria o usuário? Por que um sensor e um
atuador resolvem isso? Termine com o objetivo do projeto em uma frase.>>

\section{Pesquisa: projetos de referência}
<<Liste de 2 a 4 projetos encontrados na internet que inspiraram o de vocês. Para cada um:
o que ele faz, o que vocês aproveitaram e o que decidiram fazer diferente. Cite a fonte
na seção de Referências.>>

\section{Componentes utilizados}
\begin{table}[H]
  \centering
  \begin{tabular}{llll}
    \toprule
    \textbf{Componente} & \textbf{Função no projeto} & \textbf{Pino(s)} & \textbf{Qtde} \\
    \midrule
    Arduino Uno / ESP32     & placa de desenvolvimento & --      & 1 \\
    <<Sensor>>              & <<o que ele mede>>       & <<A0>>  & 1 \\
    <<Atuador>>             & <<o que ele faz>>        & <<9>>   & 1 \\
    <<Resistor 220 $\Omega$>> & <<limita corrente do LED>> & --  & 1 \\
    \bottomrule
  \end{tabular}
  \caption{Lista de materiais.}
\end{table}

\section{Diagrama do circuito}
<<Insira a imagem do circuito: print do Tinkercad ou esquema desenhado. Toda ligação do
código precisa aparecer aqui, com os pinos legíveis.>>

\begin{figure}[H]
  \centering
  % \includegraphics[width=0.9\textwidth]{diagrama.png}
  \fbox{\parbox[c][6cm][c]{0.9\textwidth}{\centering <<diagrama.png>>}}
  \caption{Diagrama de ligações do projeto.}
\end{figure}

\section{Funcionamento}
<<Descreva o comportamento do sistema como uma sequência: o que o sensor lê, que decisão
o código toma e o que o atuador faz. Um fluxograma simples ajuda muito.>>

\section{Código: trechos relevantes}
<<Não cole o código inteiro. Escolha os 2 ou 3 trechos que fazem o projeto ser o que é
(leitura do sensor, regra de decisão, acionamento do atuador) e explique cada um em
um parágrafo. O código completo vai no link da capa.>>

\begin{lstlisting}[caption={Leitura do sensor e decisão.}]
// <<substitua pelo seu trecho>>
void loop() {
  int leitura = analogRead(A0);
  if (leitura < LIMIAR) {
    digitalWrite(ATUADOR, HIGH);
  } else {
    digitalWrite(ATUADOR, LOW);
  }
  delay(200);
}
\end{lstlisting}

\section{Testes e ajustes}
<<O que foi testado, o que deu errado e como foi corrigido. Registre valores medidos
(limiar do sensor, tempos, distâncias). Uma tabela de testes vale mais do que um parágrafo.>>

\begin{table}[H]
  \centering
  \begin{tabular}{lll}
    \toprule
    \textbf{Teste} & \textbf{Esperado} & \textbf{Obtido} \\
    \midrule
    <<Cobrir o sensor com a mão>> & <<LED acende>> & <<OK / ajustado o limiar para 380>> \\
    \bottomrule
  \end{tabular}
  \caption{Registro de testes.}
\end{table}

\section{Resultados, limitações e melhorias}
<<O projeto cumpre o objetivo? O que ficou faltando? O que vocês fariam com mais uma semana?>>

\section{Divisão de tarefas}
\begin{table}[H]
  \centering
  \begin{tabular}{ll}
    \toprule
    \textbf{Integrante} & \textbf{Responsabilidades} \\
    \midrule
    <<Nome 1>> & <<pesquisa e proposta; montagem do circuito>> \\
    <<Nome 2>> & <<código do sensor; testes>> \\
    <<Nome 3>> & <<código do atuador; vídeo>> \\
    <<Nome 4>> & <<relatório; apresentação>> \\
    \bottomrule
  \end{tabular}
\end{table}

% =====================================================================
\section*{Referências}
\addcontentsline{toc}{section}{Referências}
\begin{itemize}
  \item <<AUTOR. Título do projeto de referência. Site, ano. Disponível em: url. Acesso em: data.>>
  \item ARDUINO. Language Reference. Disponível em: \url{https://www.arduino.cc/reference}. Acesso em: <<data>>.
  \item <<Datasheet do sensor / atuador utilizado.>>
\end{itemize}

\end{document}
